%\documentclass[preprint,12pt]{elsarticle}
\documentclass[preprint,12pt,authoryear]{elsarticle}
%\documentclass[final,3p,times,twocolumn]{elsarticle}
%\documentclass[authoryear,preprint,review,12pt]{elsarticle}
%\documentclass[final,1p,times]{elsarticle}
% \documentclass[final,1p,times,twocolumn]{elsarticle}
%\documentclass[final,3p,times]{elsarticle}
%\documentclass[final,3p,times,twocolumn]{elsarticle}
%\documentclass[final,5p,times]{elsarticle}
%\documentclass[final,5p,times,twocolumn]{elsarticle}

\usepackage{amssymb}
\usepackage{amsmath}
\usepackage{tabularx}
\usepackage{longtable}
\usepackage{tabulary}
\usepackage{array}
\usepackage{lineno}
\usepackage{natbib}
% \usepackage{algorithmicx}
\usepackage{algorithm}
\usepackage{algpseudocode}
\usepackage[singlelinecheck=off]{caption}


\journal{International Journal of Forecasting}

\begin{document}

\begin{frontmatter}

\title{Enhancing Portfolio Management: Integrating Well-Behaved Stocks with Markowitz's Mean-Variance Framework}
\author[inst1]{Elaheh Soleymanpour}
\ead{elaheh.soleymanpour@mail.um.ac.ir}
\author[inst1]{Ahmad Ghorbanipour}
\ead{ghorbanipour.ahmad@um.ac.ir}
\author[inst1]{Elham Fathipour}
\ead{elhamfathi@mail.um.ac.ir}
\author{Hadi Sadoghi Yazdi\corref{cor1}\fnref{inst1}}
\cortext[cor1]{Corresponding author}
\ead{h-sadoghi@um.ac.ir}
\affiliation[inst1]{organization={Department of Computer 
            Engineering}, addressline={Ferdowsi University of Mashhad}, city={Mashhad}, country={Iran}}

\begin{abstract}

The stock market is a vital financial arena with the potential to significantly boost a country's industrial growth through informed investment strategies. A crucial initial step in this process is analyzing optimal portfolio optimization methods. While some argue that past performance cannot predict future outcomes, a behavioral approach identifies certain stocks as "well-behaved" at specific times, suggesting that historical data can inform future trends.This article presents an innovative approach to portfolio optimization by enhancing the traditional Markowitz model with an adaptive stock selection framework using elastic matching based on Kullback-Leibler (KL) divergence. While the Markowitz method focuses on minimizing risk and maximizing returns, the proposed model identifies “well-behaved” stocks. By leveraging probability density functions derived from similar historical data windows, this approach adaptively selects stocks with consistent behavioral trends. Experimental results demonstrate that this method improves portfolio performance by an average of over 25\% compared to the basic Markowitz approach, underscoring its effectiveness in optimizing investment strategies through a behavior-informed selection process.


\end{abstract}

\begin{keyword}
Portfolio management \sep well-behaved stocks \sep Markowitz approach \sep Probability density function \sep Stock market


\end{keyword}

\end{frontmatter}

%% main text
\section{Introduction}
\label{sec:Introduction}
Today, proper capital management can be a turning point for both individual financial potential and a country's economy. If capital is directed and channeled in the right direction, it can meet the needs of both industry and society simultaneously. One of the most important capital markets is the securities exchange, through which various industrial and service companies transfer their shares to real and legal persons. After determining the investor's investment framework, the most important function is the correct arrangement of the investment portfolio. To achieve this, the best stocks should be identified and selected based on the investor's investment policies, and the amount of capital allocated to each of them should be determined. The first predictable outcome of selecting and arranging a portfolio is to maximize profitability in the shortest possible time, in other words, the rate of profitability relative to its time should be high.
A portfolio refers to the set of stocks that an investor holds in their investment basket. Portfolio optimization involves selecting $K$ out of $N$ active companies in the stock market and determining the percentage of each stock's presence in the portfolio. This process aims to allocate capital to stocks of selected companies that can lead to the highest profit. The level of portfolio risk and other characteristics of an investor can be crucial and influential in this arrangement. Two approaches are taken to solve the portfolio management problem. The first approach is an intelligent search problem that receives a large number of stocks and suggests the best ones along with the percentage of their presence in the portfolio. In the second approach, the problem is divided into two steps. Initially, $K$ stocks are selected out of $N$ stocks (where $N > K$) using certain methods, and then optimization is performed based on this selection.
Markowitz's method is one of the most significant and fundamental methods used in portfolio management modeling. In 1952, he presented a model based on the  {\ttfamily"}Expected Return-Variance{\ttfamily"} law, which has since become the basis of many studies( \cite{markowitz1952}). In his article, Markowitz argues that the expected return should be maximized, and the risk of the stock portfolio should be minimized. \ref{Markowitz} presents the Markowitz-based objective function with the parameter $\lambda$, based on \cite{markowitz1959portfolio}:

\begin{equation}\label{Markowitz}
\begin{split}
& min \lambda\left[ \sum_{i=1}^{N}\sum_{j=1}^{N}w_i w_j\sigma_{ij} \right] - (1-\lambda)\left[ \sum_{i=1}^{N}w_i\mu_i \right] \\
& s.t.\begin {cases}
    \sum_{i=1}^{N} w_i=1 & \\ \sum_{i=1}^{N} z_i=K & \\ \varepsilon_i z_i\le w_i\le \sigma_i z_i & i=1, ..., N \\ z_{i} \in {0,1} & i=1, ..., N \\ 0\le w_i\le 1 & i=1, ..., N \\ 0\le \varepsilon_{i}\le \sigma_{i}\le 1 & i=1, ..., N 
\end {cases}
\end{split}
\end{equation}

The first and second terms in the above objective function represent risk and portfolio return optimization, respectively. The variable $w_{i}$ represents the weight of stock $i$ in the portfolio, $\mu_i$ is the return of stock $i$, and $\sigma_{ij}$ is the $(i,j)th$ element of the covariance matrix. $\epsilon_{i}$ and $\delta_i$ define the minimum and maximum percentage of stock $i$, respectively. In this approach, $z_i$  indicates the presence or absence of stock $i$ in the portfolio, which can be initialized in the optimization or in a pre-optimization step.

Recent advancements in portfolio optimization reflect a broad array of improvements to the traditional Markowitz model, enhancing its robustness and extending its applications across diverse sectors. A key area of progress involves applying Markowitz analysis to new domains, such as sustainable energy. \cite{Botero2024} applied the model to Colombia's hydro-dominated energy sector, finding that integrating wind power reduces risk and price volatility, thereby improving the efficient frontier. This adaptation highlights the potential of Markowitz-based approaches in sustainable energy planning. Similarly, addressing limitations of the classic mean-variance model,  \cite{Ortiz2023} employed covariance matrix shrinkage with ridge regression to stabilize portfolio weights and improve out-of-sample Sharpe ratios, thus making the model more practical. In the context of urban power decarbonization, \cite{Yang2024} integrated deep learning and green hydrogen storage into the Markowitz framework to enhance returns amid the challenges of renewable energy.\cite{Hui2024} proposed a Factor-Augmented Derating (FAD) technique for large-scale portfolios, narrowing the gap between expected and realized returns, with significant implications for stock market applications. \cite{Petukhina2024} further improved model stability by using projected gradient descent, which reduces sensitivity to input fluctuations and transaction costs, bolstering portfolio stability and out-of-sample performance.

Building on these core advancements, various innovative methods are being developed to refine stock selection, risk management, and overall portfolio performance. \cite{cheong2017using} introduced a two-stage approach combining k-means clustering for stock selection with genetic algorithms for optimizing portfolio weights. \cite{farid2022portfolio} employed self-organizing maps to cluster companies, followed by genetic algorithms to optimize selected stocks, creating a well-balanced portfolio. To address market instability, \cite{pahade2021multi} introduced a fuzzy logic model that considers variance and semi-variance, capturing market volatility for more robust optimization. Other structured search techniques have also emerged; for instance, \cite{akbay2020parallel} used neighborhood search algorithms with quadratic programming (QP) for optimal stock weights, while \cite{kizys2019biased} integrated assignment problem solutions with repetitive local search and QP for enhanced accuracy.

Sparse and short-term-focused models further contribute to diversified strategies. \cite{lai2018short} proposed a short-term, sparse model that simplifies portfolio composition by emphasizing the immediate profitability of selected stocks. Bayesian learning is another area of innovation, where \cite{de2019bayesian} used Bayesian algorithms with Gaussian assumptions to refine stock return estimates through prior and posterior distributions. For multi-period modeling,\cite{li2022multi} demonstrated that stepwise adjustments to portfolios over multiple periods could yield higher returns compared to single-period models. \cite{mercurio2020entropy} introduced an entropy minimization approach as an alternative to variance-based risk measures, creating a diversified portfolio structure that balances return and risk.

Advanced predictive and time-varying models have also gained traction. (\cite{katsikis2022time}) explored time-varying mean-variance datasets with time-variant quadratic programming (TVQP), while(\cite{kaczmarek2022building}) utilized machine learning to dynamically predict stock performance and optimize allocations. Meanwhile, (\cite{torrente2022rescaling}) applied game-theoretic adjustments to classical models, introducing competitive rebalancing for portfolio stability. (\cite{arcuri2023robust}) employed a Value at Risk (VaR) model that emphasizes strong conditional VaR to maintain stable, long-term performance.

Additional perspectives further enhance portfolio optimization, such as combining fundamental indexing (FI) with mean-variance optimization (\cite{pysarenko2019predictive}), and developing multivariate affine generalized hyperbolic (MAGH) frameworks (\cite{chaweewanchon2022markowitz}) for asset allocation. High-dimensional data in portfolio optimization is addressed by \cite{li2022spectrally}, who proposed a method for handling large, high-dimensional datasets effectively. \cite{katsikis2022time} introduced time-varying quadratic programming (TVQP) to manage datasets with changing means and variances over time.


This article introduce a novel concept called {\ttfamily"}well-behavedness of stocks{\ttfamily"} based on a retrospective idea.This approach predicts future stock performance by analyzing probability density functions of similar historical windows, thereby identifying stocks with consistent past behaviors likely to inform future trends. Ultimately, a suitable portfolio selection approach is combined with the main Markowitz approach. The algorithm's effectiveness is evaluated on the Tehran Stock Exchange and other stock exchanges worldwide. Section \ref{sec:Experiments} explains the proposed methods in detail in several subsections. In section \ref{sec:Discussion}, the authors introduce the dataset used, evaluation criteria, and the results of applying the proposed method and selecting effective portfolios. The final section concludes with a discussion and conclusion.

\section{Proposed Method}
\label{sec:Proposed}


The proposed method builds on the traditional Markowitz approach (\cite{markowitz1959portfolio}) and incorporates an adaptive framework for stock selection using elastic matching based on Kullback-Leibler (KL) divergence. The overall aim is to enhance portfolio optimization by identifying well-behaved stocks whose past behaviors can reliably predict future performance. The method is divided into several steps, as described below:


\subsection{Identifying Best Similar Window based on the function of similarity degree(FBS)}\label{sec:Find} 

Let the price of a stock over a time period $( T )$ be represented as a time series \( X_T = \{x_t\}_{t=1}^T \), where \( x_t \) is the price at time \( t \). The return signal \( r_t \) is derived from the price series, and it is used in further calculations for portfolio construction.

To predict future stock behavior, the proposed method searches for a similar window in the past. Assuming that the current window of interest starts at time \( t \) with a length of \( w \) days, the algorithm compares this window to past windows of the same stock. However, instead of searching the entire history, the search is limited to \( w'' \) days prior to \( t \).  the pseudo-code for the similarity search is provided in Algorithm \ref{alg:best-match}.

\begin{algorithm}
\caption{Finding the best matching window in stock history based on similarity degree}\label{alg:best-match}
\begin{algorithmic}
\State $ Best-similar-index \gets t$
\State $Best-divergence \gets 1000$
\For{$t^{\prime}=t+w:t+w^{\prime}$}
\State $divergence \gets D(r(t:t+w)\parallel (r:t^{\prime}:t^{\prime}+w))$
\If{$ divergence < Best-divergence$}
\State $Best-divergence \gets divergence$
\State $Best-similar-index \gets t^{\prime}$
\EndIf
\EndFor
\State \textbf{return} $Best-similar-index, Best-divergence$
\end{algorithmic}
\end{algorithm}
The divergence \( D \) between two windows is measured using kernel density estimates (KDE), as outlined below.

\subsection{Calculating the Degree of Divergence}
\label{sec:Function}

The divergence or dissimilarity between two windows is calculated using their probability density functions (PDFs). Instead of treating each window as a single unit, the method divides each \( w \)-day window into overlapping sub-windows of length \( w' \). 

For each sub-window pair \( i \), the divergence is calculated using a kernel density estimation (KDE), defined by :

\begin{equation} \label{kernel}
    \hat{f}_{t_i} = \frac{1}{n} \sum_{j=t+i}^{t+i+w'} K_h(x - r_j)
\end{equation}

Where \( K_h \) is the Gaussian kernel function. The total divergence is computed by averaging the divergences across all sub-windows by \ref{Div} :

\begin{equation} \label{Div}
    \text{div} = \sum_{i=0}^{w-w'}D(\hat{f}_{t_i} \parallel \hat{f}_{t'_i})
\end{equation}

This process is illustrated in Algorithm  \ref{alg:kde-divergence}, which provides the steps for calculating the dissimilarity between two windows.

\begin{algorithm}
\caption{Pseudocode for calculating the dissimilarity between two windows relative to each other}
\label{alg:kde-divergence}
\begin{algorithmic}
\State $div \gets 0$
\For{$i = 0$ \textbf{to} $w-w^{\prime}$}
\State $\text{Find}\; f1 = \text{KDE of}\; r(t+i:t+i+w^{\prime})$
\State $\text{Find}\; f2 = \text{KDE of}\; r(t^{\prime}+i:t^{\prime}+i+w^{\prime})$
\State $div \gets$ divergence of $f1$ and $f2$
\State $div \gets div + div$
\EndFor
\State \textbf{return} $div$
\end{algorithmic}
\end{algorithm}


where $\hat{f}_{t_i}$ and $\hat{f}_{t_j}$ are kernel density functions of two distributions  \ref{phat} and \ref{pprimehat} explain how these functions are calculated based on the Gaussian kernel:

\begin{equation} \label{phat}
    \hat{f}_{t_i}= \frac{1}{n(2\pi \sigma^2)^\frac{d}{2}}\sum_{i=1}^{n}exp(-\frac{\left\| x-x_{i} \right\|^2}{2\sigma ^2})
\end{equation}

\begin{equation} \label{pprimehat}
    \hat{f}_{t_j}= \frac{1}{n^{\prime}(2\pi {\sigma^{\prime}}^2)^\frac{d}{2}}\sum_{i^{\prime}=1}^{n^{\prime}}exp(-\frac{\left\| x-x^{\prime}_{i^{\prime}} \right\|^2}{2{\sigma^{\prime}}^2 })
\end{equation}

If the Jensen-Shannon divergence measure is used to calculate the distance between these two distributions, it should be calculated according to:
\begin{equation}\label{JS}
    D_{JS}(\hat{f}_{t_i}\parallel \hat{f}_{t_j})=\frac{1}{2}\left(D_{KL}\left[ \hat{f}_{t_i}\parallel\frac{\hat{f}_{t_i}+\hat{f}_{t_j}}{2} \right] + D_{KL}\left[ \hat{f}_{t_j}\parallel\frac{\hat{f}_{t_i}+\hat{f}_{t_j}}{2} \right]\right)
\end{equation}

KL stands for Kullback-Leibler divergence between two sub-windows \( t_i \) and \( t_j \) is defined as:

\begin{equation}\label{eq:kl_divergence_subwindows}
    D_{\text{KL}}(\hat{f}_{t_i} \parallel \hat{f}_{t_j}) = \sum_{x \in \mathcal{X}} \hat{f}_{t_i}(x) \log\left(\frac{\hat{f}_{t_i}(x)}{\hat{f}_{t_j}(x)}\right),
\end{equation}

where \( \mathcal{X} \) is the set of possible stock prices.


\subsection{Elastic Matching and Well-Behaved Stocks}
\label{sec:welbehaved}


The elastic matching condition evaluates the similarity between the recent behavior of a stock's price distribution and its past behavior by computing the mean Kullback-Leibler (KL) divergence. Let \( T_1 \) represent the most recent data window and \( T_i \) for \( i = 2, \dots, n \) represent past data windows. The mean KL divergence \( \overline{D_{KL}} \) between the current window \( T_1 \) and the past windows \( T_i \) is given by:

\begin{equation}\label{eq:D_KL}
    \overline{D_{KL}} = \frac{1}{n-1} \sum_{i=2}^{n} D_{KL}(F_{T_1} \parallel F_{T_i}).
\end{equation}

If the mean divergence \( \overline{D_{KL}} \) is smaller than a predefined threshold \( \epsilon \), i.e., \( \overline{D_{KL}} < \epsilon \), the stock's behavior over these sub-windows is considered elastic. This means there is a consistent pattern between the present and past distributions.

A stock is considered well-behaved if the elastic matching condition continues to hold over an extended period \( T + \Delta T \), where \( n' = \lfloor \Delta T / \tau \rfloor \) future sub-windows \( T_i \) for \( i = n+1, \dots, n+n' \) are examined. The stock's price distribution must satisfy:

\begin{equation}\label{eq:wel}
    \frac{1}{n'} \sum_{i=n+1}^{n+n'} D_{KL}(F_{T_1} \parallel F_{T_i}) \leq \epsilon.
\end{equation}

When this condition is met, the stock is interpreted as having predictable future behavior based on its past, indicating good performance over time.

\begin{algorithm}
\caption{The pseudocode to find the top 35 well-behaved stocks}
\label{alg:pseudocode}
\begin{algorithmic}
\State $\text{Indexs} \gets \{\}$
\For{$i=1$ \textbf{to} $N$}
\State $numberOfTrueCondition \gets 0$
\State $totalNumber \gets 0$
\For{$t=21$ \textbf{to} $\text{length}(r_i)$}
\State $\mu_{it} \gets \text{mean}(r_i(t-w:t))$
\State $\text{bestSimilarIndex}, \text{bestSimilarity} \gets \text{FBS}(r_i,t)$
\State $\mu_{it}' \gets \text{mean}(r_i(\text{bestSimilarIndex}-w:\text{bestSimilarIndex}))$
\If{$\text{bestSimilarity} < 1$}
\State $totalNumber \gets totalNumber + 1$
\If{$\mu_{it}/\mu_{it}' > 0$}
\State $numberOfTrueCondition \gets numberOfTrueCondition + 1$
\EndIf
\EndIf
\EndFor
\State $\text{Indexs(end+1)}.\text{index} \gets i$
\State $\text{Indexs(end)}.\text{percentage} \gets numberOfTrueCondition/totalNumber$
\EndFor
\State $\text{sortedIndex} \gets \text{sort}(\text{Indexs}, -\text{percentage})$
\State \textbf{return} $\text{sortedIndex}(1 : 35)$
\end{algorithmic}
\end{algorithm}

Algorithm\ref{alg:pseudocode} outlines the pseudocode used to identify well-behaved stocks. For each stock \( i \), we iterate over its price data \( r_i \), calculate the average return \( \mu_{it} \) over a window of length \( w \), and find the most similar past window using the FBS method. If the similarity score is below a threshold (typically \( <1 \)), and the ratio \( \mu_{it} / \mu'_{it} > 0 \) (where \( \mu'_{it} \) represents the average return of the similar window), the stock is classified as well-behaved for that window. The percentage of such well-behaved instances is computed over the entire time series, and the top 35 stocks with the highest percentage are selected for portfolio management.

\subsection{integrating Well-Behaved Stocks with the Markowitz's basic process}
\label{sec:combining}


Once the top well-behaved stocks are identified, the FBS method is integrated into the Markowitz optimization framework. In this approach, two data windows are considered: the current window and a future window identified as similar to the current one.As shown in Fig.\ref{fig:similarwindow}, the black window represents the current data window, while the green window represents the unknown future. The red window represents a past window that is similar to the current one, identified using the FBS algorithm. The purple window represents the future of this similar window. If the future of the similar window closely resembles the current window's future, this stock behavior is considered a good basis for portfolio decision-making.

\begin{figure}
\begin{center}
  \includegraphics[width=0.99\linewidth]{similarwindow.jpg}
  \caption{The meaning of current, future, similar, and future  of similar window in an example of implementing the FBS algorithm for a specific stock}
  \label{fig:similarwindow}  
\end{center}
\end{figure}

To quantify the similarity between the current and similar windows, we calculate the Jensen-Shannon (JS) divergence over 10 windows. If the total divergence is less than 1, the performance is typically better. A smaller divergence value indicates higher similarity, and if the divergence between the two windows is less than 0.5, more emphasis is placed on the future of the similar window for prediction.

The weight coefficient \( \gamma_i \) for the \( i \)-th stock in the portfolio is defined as follows:

\begin{equation}\label{gamma}
    \gamma_i = 1 - \min\left(1, \frac{0.5}{d}\right),
\end{equation}

where \( d \) is the divergence between the current and similar windows. This coefficient adjusts the importance placed on the future window versus the current window.

The Markowitz optimization model is then modified to account for both windows. The objective function becomes:


\begin{equation}\label{min}
\begin{split}
    min(& \lambda \left\lceil \sum_{i=1}^{N} \gamma_i \sum_{j=1}^{N}w_i w_j\sigma_{ij} \right\rceil
    -(1-\lambda) \left[ \sum_{i=1}^{N}\gamma_i w_i \mu_i \right]) \\ & + (\lambda^{\prime} \left\lceil \sum_{i=1}^{N} (1-\gamma_i) \sum_{j=1}^{N}w_i w_j \sigma^{\prime}_{ij} \right\rceil - (1-\lambda^{\prime}) \left[ \sum_{i=1}^{N}(1-\gamma_i) w_i \mu^{\prime}_i \right])
\end{split}
\end{equation}

where \( \lambda \) and \( \lambda' \) represent risk coefficients, \( \sigma_{ij} \) and \( \sigma'_{ij} \) are the covariances of the stock returns in the current and future windows, and \( \mu_i \) and \( \mu'_i \) are the mean returns for the current and future windows, respectively.

By solving this optimization problem using quadratic programming, the desired weight allocation for each stock in the portfolio is determined. The weight constraints ensure that each stock holds between 5\% and 40\% of the portfolio. This method ensures that the selected stocks not only show consistent behavior in the past but are also positioned to perform well in the future, improving the overall portfolio's stability and return potential.




\section{Experiments}\label{sec:Experiments}
In the proposed method, several key parameters are set to optimize performance. The window length \(w\) is set to 20 days for both the proposed method and the Markowitz process. The rolling window length \(w′\), used for comparing different windows, is set to 10 days, resulting in smaller 10-day windows within each 20-day window. Additionally, the parameter \(w′′\) determines the look-back period for identifying similar windows and is set to 500 days (roughly two years). To calculate the appropriate bandwidth for kernel density estimation (KDE), the improved Sheater and Jones method is applied to the daily return signals of each stock. This approach is consistently used for all stages involving the stock’s performance evaluation.

\subsection{Experimental data} \label{sec:Experimental data}
The experiments were conducted on daily stock price data from the Tehran Stock Exchange. Stocks that did not meet the behavior criteria, such as {\ttfamily"}Exchange-Traded Fund{\ttfamily"} and {\ttfamily"}Rights Offering Shares{\ttfamily"}, were excluded. Ultimately, 368 symbols were selected for testing the proposed algorithm. This selection ensures that only stocks exhibiting favorable characteristics for exploration are included in the analysis.


\subsection{Evaluation by Back Testing Method}
\label{sec:Eval}
To evaluate the performance of the proposed method, a back-testing system was implemented. Based on the system’s suggested coefficients for portfolio management at time \(t\), stocks were purchased. The algorithm measures the average daily return over an \(N\)-day window, using multiple-day windows for buying and selling decisions. At the end of each window, profits and losses are calculated, and the process is repeated over a total period of 1000 days. The evaluation is reported in ten 100-day intervals, with both short-term and overall average returns used to assess the algorithm’s effectiveness. A better performance over time indicates a more successful portfolio management strategy.

\subsection{FBS Method and Bandwidth Selection}
\label{sec:Eval FBS}

The FBS (Find Best Similarity) algorithm was tested to determine the optimal window size. Fig.\ref{fig:fbs} demonstrates the results for the “Foolad” stock symbol on day 537, where the red window represents the current time and the green window represents a similar past window. The future trends of both windows are closely aligned, showcasing the FBS algorithm's capability to identify key patterns. Additionally, bandwidth selection for KDE was tested using different methods. The improved Sheater and Jones method consistently outperformed other approaches. Fig.\ref{fig:comparing} illustrates how this method provided superior accuracy in predicting future trends, leading to its selection for the final approach.


\begin{figure}
\begin{center}
    \includegraphics{fbs.jpg}
    \caption{Executing the proposed FBS algorithm on the {\ttfamily"}Foolad{\ttfamily"} stock symbol, on the 537th day from the end}
    \label{fig:fbs}
\end{center}
\end{figure}




\begin{figure}
\begin{center}
  \includegraphics[width=0.9\linewidth]{comparing.jpg}
  \caption{Comparing different bandwidths for assessing the accuracy of the first condition using the best matching window when the similarity value exceeds a specified threshold}  
  \label{fig:comparing}
\end{center}
\end{figure}

\subsection{Identifying Well-Behaved Stocks}

The method defines well-behaved stocks based on ranking. Initially, stocks from the last 1000 days were analyzed, but to avoid bias from known back-testing results, this period was excluded from the final analysis. Instead, we focused on identifying well-behaved stocks from the preceding 2000 days, ensuring a forward-looking approach. Table \ref{table:finding35} presents the top 35 well-behaved symbols from the second thousand-day period, which were then used to continue the evaluation process for the first thousand days, maintaining the integrity of the test results.



\begin{longtable}{ p{.05\textwidth} p{.13\textwidth} p{.14\textwidth} p{.25\textwidth} p{.13\textwidth} p{.15\textwidth} } 
\caption{Lists the top 35 stocks identified using the Sheater and Jones bandwidth method, considering only the best similar windows that exceed a specified similarity threshold, while excluding the last thousand days of data.}
\label{table:finding35}
\\ \hline 
\textbf{Row} & \textbf{Symbol Stock Name} & \textbf{Accuracy Percentage}& \textbf{Industry} & \textbf{Total Number of Data Point} & \textbf{Number of instances where divergence exceeds the threshold} \\ \hline
\endhead
1 & Zegoldasht & 100 & Farm and related services & 125 & 1 \\ 
3 & Ghgolsta & 81.8182 & Food and Beverage Products & 125 & 9 \\ 
4 & Koravi & 81.8182 & Metal ore extraction & 444 & 9 \\ 
5 & Benirou & 76.4706 & Electrical machinery and equipment & 397 & 52 \\ 
6 & Setran & 74.1379 & Cement, lime, and gypsum & 454 & 43 \\ 
7 & Seydkoo & 72.0588 & Cement, lime, and gypsum & 166 & 49 \\ 
8 & VaAzar & 69.2308 & Construction, real estate, and properties & 363 & 54 \\ 
9 & Sheklar & 68.75 & Chemical products & 367 & 11 \\ 
10 & Kama & 68.2927 & Metal ore extraction & 376 & 28 \\ 
11 & Sekhrz & 68.1034 & Cement, lime, and gypsum & 379 & 79 \\ 
12 & FarAvar & 67.8571 & Basic metals & 383 & 19 \\ 
13 & Sepaha & 67.3469 & Cement, lime, and gypsum & 436 & 33 \\ 
14 & Chekaveh & 66.6667 & Paper products & 400 & 32 \\ 
15 & Shenfat & 65.6051 & Petroleum, coke, and nuclear fuel products & 451 & 103 \\ 
16 & Ghadasht & 65.285 & Food and Beverage Products excluding sugar and sweets & 376 & 126 \\ 
17 & Nbours & 64.7059 & Auxiliary activities to financial intermediaries & 117 & 11 \\ 
18 & Kenour & 63.2911 & Metal ore extraction & 177 & 50 \\ 
19 & Barkat & 62.5 & Pharmaceutical materials and products & 114 & 10 \\ 
20 & Kesaveh & 62.2642 & Tiles and ceramics & 372 & 33 \\ 
21 & Kesaedi & 61.8421 & Tiles and ceramics & 360 & 47 \\ 
22 & Deshimi & 61.7647 & Pharmaceutical materials and products & 377 & 84 \\ 
23 & Faspa & 61.5 & Basic metals & 362 & 123 \\ 
24 & Fasmin & 61.2903 & Basic metals & 410 & 19 \\ 
25 & Kebourse & 61.1111 & Auxiliary activities to financial intermediaries & 117 & 11 \\ 
26 & Bourse & 61.1111 & Auxiliary activities to financial intermediaries & 117 & 11 \\ 
27 & Ghapira & 61.039 & Sugar and sweets & 375 & 47 \\ 
28 & Maroun2 & 60.5839 & Chemical products & 238 & 83 \\ 
29 & Sarbil & 60.2804 & Cement, lime, and gypsum & 325 & 129 \\ 
30 & Maroun & 59.4203 & Chemical products & 238 & 82 \\ 
31 & Fakhas & 59.0517 & Basic metals & 331 & 137 \\ 
32 & Sashahed & 58.9744 & Construction,real estate, and properties & 365 & 23 \\ 
33 & Depars & 58.9286 & Pharmaceutical materials and products & 458 & 99 \\ 
34 & Safha & 58.7065 & Cement, lime, and gypsum & 346 & 118 \\ 
35 & Fama & 58.6319 & Metal products manufacturing & 405 & 180 \\ \hline
% needs to go inside longtable environment
\label{tab:myfirstlongtable}
\end{longtable}


\subsection{Comparison of the Proposed Method and the Markowitz Baseline Method}
In this section, we utilize the Markowitz baseline method to calculate the mean return rate for the 50 stock portfolios obtained in the previous section, employing a back-testing approach. We repeat this process for the proposed method and then compare the results from both methods. By subtracting the daily average return of the Markowitz method from that of the proposed method, we can draw conclusions regarding the superiority of one approach over the other and assess the degree of improvement achieved as shown in Table \ref{table:improvement}.

\begin{longtable}{ p{.15\textwidth}  p{.27\textwidth} p{.25\textwidth} p{.25\textwidth} }
\caption{The percentage of improvement in the proposed method compared to the Markowitz baseline method for 50 different portfolios.}
\label{table:improvement}
% \begin{tabularx}{\textwidth}{|c|X|X|X|}
\hline 
\textbf{Portfolio} & \textbf{The result of subtracting B from A, divided by the absolute value of B, multiplied by 100} & \textbf{Average return using the Markowitz baseline method (B)} & \textbf{Average return using the proposed method (A)} \\ 
\hline
\endhead
1& 20.7415 & 0.2573 & 0.3106 \\
2 & 56.2041 & 0.1801 & 0.2814 \\
3 & 45.4560 & 0.2566& 0.3732 \\
4 & 152.7466& 0.1062& 0.2683 \\
5& 15.6404& 0.1904& 0.2202 \\
6& 68.2275& 0.1604& 0.2698 \\
7 & 57.6507& 0.1643& 0.2591 \\
8 & 63.2526& 0.2035& 0.3322                    \\
9& 1.4891& 0.2911& 0.2955                    \\
10& 24.5071& 0.2753& 0.3428                    \\
11 & 12.344 & 0.2757 & 0.2417                    \\
12 & 104.9510 & 0.1340& 0.2747                    \\
13 & 57.7825& 0.1292& 0.2039                    \\
14 & 8.2209& 0.3380& 0.3103                    \\
15 & 2.3804 & 0.2232& 0.2285                    \\
16 & 9.7994 & 0.2241& 0.2460                    \\
17 & 38.6751 & 0.1954& 0.2709                    \\
18 & 7.1199 & 0.2488 & 0.2665                    \\
19 & 45.8413& 0.1851 & 0.2700                    \\
20 & 1.9877& 0.2048 & 0.2088                    \\
21 & 40.2888 & 0.2051 & 0.2878                    \\
22 & 17.1097  & 0.1616 & 0.1893                    \\
23 & 50.8998 & 0.2845  & 0.4293                    \\
24 & 16.3096  & 0.2214 & 0.2575                    \\
25 & 42.4601  & 0.2629 & 0.1513                    \\
26 & 26.5798  & 0.2824  & 0.3575                    \\
27 & 15.8851 & 0.2045 & 0.2369                    \\
28 & 46.1852 & 0.1982 & 0.2897  \\
29 & 49.9436 & 0.2034 & 0.3050                    \\
30 & 3.7192 & 0.2909 & 0.2801                    \\
31 & 4.0503 & 0.2324  & 0.2230                    \\
32 & 2.3509 & 0.2526 & 0.2585                    \\
33 & 34.7096  & 0.1422 & 0.1915                    \\
34 & 9.3057  & 0.2628  & 0.2873                    \\
35 & 8.2149  & 0.2700 & 0.2922                    \\
36 & 44.9178 & 0.2109 & 0.3057                    \\
37 & 54.6303 & 0.2706 & 0.4184                    \\
38 & 17.3260 & 0.2983 & 0.3500                    \\
39 & 19.3977 & 0.2257 & 0.2695                    \\
40 & 9.2423 & 0.1750 & 0.1912                    \\
41 & 34.1795 & 0.2613 & 0.3505                    \\
42 & 39.0505 & 0.2474 & 0.3440                    \\
43 & 26.8253 & 0.2657 & 0.1944                    \\
44 & 15.5447 & 0.1981 & 0.1673                    \\
45 & 127.3386  & 0.1286 & 0.2923                    \\
46 & 45.3645 & 0.1488 & 0.2163                    \\
47 & 12.2787 & 0.2471 & 0.2774                    \\
48  & 22.2219 & 0.2218 & 0.2711                    \\
49 & 47.8972 & 0.2152 & 0.3183                    \\
50 & 32.5290 & 0.2455 & 0.3254                    \\ \hline
% \end{tabularx}
\noalign{\smallskip}
\end{longtable}


\section{Discussion and Conclusion}\label{sec:Discussion}
Table \ref{table:comparing}  the performance improvements of the proposed method compared to the baseline Markowitz approach across fifty portfolios. The behavior of the portfolios was evaluated during the second thousand days, while the back-testing was conducted over the most recent thousand days. The results clearly highlight the benefits of the proposed method in identifying well-behaved stocks and leveraging future windows to enhance performance.

\begin{table}
\caption{Comparison of the proposed method with the Markowitz baseline for 50 portfolios of well-behaved stocks over the last 1000 days.}
\label{table:comparing}
\begin{tabular}{c  p{0.7\linewidth}}
  \hline
  \textbf{Proposed method} & \multicolumn{1}{c}{\textbf{Results}} \\ 
  \hline
  7\% & Percentage of portfolios with negative improvement, indicating worse performance compared to the Markowitz baseline.  \\ 
  
  86\% & Success rate of the proposed method, calculated as the number of portfolios with improved performance divided by the total number of portfolios. $(\frac{43}{50}*100)$ \\ 
  
  29.7\% & Average percentage improvement across all cases. \\
  
  37.2\% & Average percentage improvement in cases where the method led to better performance. \\
  \hline
\end{tabular}
\label{tab:compared}
\end{table}

The result shows that, when well-behaved stocks are accurately identified and the proposed method is applied, significant improvements in portfolio performance can be achieved. This reinforces the value of integrating the current and future windows in the optimization process for more accurate forecasting and enhanced portfolio returns.


Future work can improve upon the current methodology by introducing intelligent stock selection criteria, such as using fundamental stock indicators or clustering based on recent performance, rather than random selection. Additionally, the Markowitz optimization process could be enhanced by incorporating variable or multi-length windows and adjusting the risk-taking coefficient \( \lambda \) over time to better capture market dynamics.

Other potential enhancements include refining the back-testing process by accounting for trading factors like buy and sell queues, which would more accurately reflect real-world conditions. Improvements to the FBS algorithm, such as better distance learning between similar future windows and their actual outcomes, along with the use of variable bandwidth over time, could also lead to more precise and adaptive forecasting models.

%% If you have bibdatabase file and want bibtex to generate the
%% bibitems, please use
%%
 %\bibliographystyle{plainnat} 
 %\bibliographystyle{agsm} % Name-Date (Harvard) style
 %\bibliography{cas-refs}



\bibliographystyle{elsarticle-harv} 
\bibliography{cas-refs}

\end{document}
\endinput
%\documentclass[preprint,12pt]{elsarticle}
\documentclass[preprint,12pt,authoryear]{elsarticle}
%\documentclass[final,3p,times,twocolumn]{elsarticle}
%\documentclass[authoryear,preprint,review,12pt]{elsarticle}
%\documentclass[final,1p,times]{elsarticle}
% \documentclass[final,1p,times,twocolumn]{elsarticle}
%\documentclass[final,3p,times]{elsarticle}
%\documentclass[final,3p,times,twocolumn]{elsarticle}
%\documentclass[final,5p,times]{elsarticle}
%\documentclass[final,5p,times,twocolumn]{elsarticle}

\usepackage{amssymb}
\usepackage{amsmath}
\usepackage{tabularx}
\usepackage{longtable}
\usepackage{tabulary}
\usepackage{array}
\usepackage{lineno}
\usepackage{natbib}
% \usepackage{algorithmicx}
\usepackage{algorithm}
\usepackage{algpseudocode}
\usepackage[singlelinecheck=off]{caption}


\journal{Knowledge-Based Systems }

\begin{document}

\begin{frontmatter}

\title{Enhancing Portfolio Management: Integrating Well-Behaved Stocks with Markowitz's Mean-Variance Framework}

\author[inst1]{Ahmad Ghorbanipour}
\ead{ghorbanipour.ahmad@um.ac.ir}
\author[inst1]{Elaheh Soleymanpour}
\ead{elaheh.soleymanpour@gmail.com}
\author{Hadi Sadoghi Yazdi\corref{cor1}\fnref{inst1}}
\cortext[cor1]{Corresponding author}
\ead{h-sadoghi@um.ac.ir}
\affiliation[inst1]{organization={Department of Computer 
            Engineering}, addressline={Ferdowsi University of Mashhad}, city={Mashhad}, country={Iran}}

\begin{abstract}

The stock market is a vital financial arena with the potential to significantly boost a country's industrial growth through informed investment strategies. A crucial initial step in this process is analyzing optimal portfolio optimization methods. While some argue that past performance cannot predict future outcomes, a behavioral approach identifies certain stocks as "well-behaved" at specific times, suggesting that historical data can inform future trends.This article presents an innovative approach to portfolio optimization by enhancing the traditional Markowitz model with an adaptive stock selection framework using elastic matching based on Kullback-Leibler (KL) divergence. While the Markowitz method focuses on minimizing risk and maximizing returns, the proposed model identifies “well-behaved” stocks. By leveraging probability density functions derived from similar historical data windows, this approach adaptively selects stocks with consistent behavioral trends. Experimental results demonstrate that this method improves portfolio performance by an average of over 25\% compared to the basic Markowitz approach, underscoring its effectiveness in optimizing investment strategies through a behavior-informed selection process.


\end{abstract}

\begin{keyword}
Portfolio management \sep well-behaved stocks \sep Markowitz approach \sep Probability density function \sep Stock market


\end{keyword}

\end{frontmatter}

%% main text
\section{Introduction}
\label{sec:Introduction}
Today, proper capital management can be a turning point for both individual financial potential and a country's economy. If capital is directed and channeled in the right direction, it can meet the needs of both industry and society simultaneously. One of the most important capital markets is the securities exchange, through which various industrial and service companies transfer their shares to real and legal persons. After determining the investor's investment framework, the most important function is the correct arrangement of the investment portfolio. To achieve this, the best stocks should be identified and selected based on the investor's investment policies, and the amount of capital allocated to each of them should be determined. The first predictable outcome of selecting and arranging a portfolio is to maximize profitability in the shortest possible time, in other words, the rate of profitability relative to its time should be high.
A portfolio refers to the set of stocks that an investor holds in their investment basket. Portfolio optimization involves selecting $K$ out of $N$ active companies in the stock market and determining the percentage of each stock's presence in the portfolio. This process aims to allocate capital to stocks of selected companies that can lead to the highest profit. The level of portfolio risk and other characteristics of an investor can be crucial and influential in this arrangement. Two approaches are taken to solve the portfolio management problem. The first approach is an intelligent search problem that receives a large number of stocks and suggests the best ones along with the percentage of their presence in the portfolio. In the second approach, the problem is divided into two steps. Initially, $K$ stocks are selected out of $N$ stocks (where $N > K$) using certain methods, and then optimization is performed based on this selection.
Markowitz's method is one of the most significant and fundamental methods used in portfolio management modeling. In 1952, he presented a model based on the  {\ttfamily"}Expected Return-Variance{\ttfamily"} law, which has since become the basis of many studies( \cite{markowitz1952}). In his article, Markowitz argues that the expected return should be maximized, and the risk of the stock portfolio should be minimized. \ref{Markowitz} presents the Markowitz-based objective function with the parameter $\lambda$, based on \cite{markowitz1959portfolio}:

\begin{equation}\label{Markowitz}
\begin{split}
& min \lambda\left[ \sum_{i=1}^{N}\sum_{j=1}^{N}w_i w_j\sigma_{ij} \right] - (1-\lambda)\left[ \sum_{i=1}^{N}w_i\mu_i \right] \\
& s.t.\begin {cases}
    \sum_{i=1}^{N} w_i=1 & \\ \sum_{i=1}^{N} z_i=K & \\ \varepsilon_i z_i\le w_i\le \sigma_i z_i & i=1, ..., N \\ z_{i} \in {0,1} & i=1, ..., N \\ 0\le w_i\le 1 & i=1, ..., N \\ 0\le \varepsilon_{i}\le \sigma_{i}\le 1 & i=1, ..., N 
\end {cases}
\end{split}
\end{equation}

The first and second terms in the above objective function represent risk and portfolio return optimization, respectively. The variable $w_{i}$ represents the weight of stock $i$ in the portfolio, $\mu_i$ is the return of stock $i$, and $\sigma_{ij}$ is the $(i,j)th$ element of the covariance matrix. $\epsilon_{i}$ and $\delta_i$ define the minimum and maximum percentage of stock $i$, respectively. In this approach, $z_i$  indicates the presence or absence of stock $i$ in the portfolio, which can be initialized in the optimization or in a pre-optimization step.

Recent advancements in portfolio optimization reflect a broad array of improvements to the traditional Markowitz model, enhancing its robustness and extending its applications across diverse sectors. A key area of progress involves applying Markowitz analysis to new domains, such as sustainable energy. \cite{Botero2024} applied the model to Colombia's hydro-dominated energy sector, finding that integrating wind power reduces risk and price volatility, thereby improving the efficient frontier. This adaptation highlights the potential of Markowitz-based approaches in sustainable energy planning. Similarly, addressing limitations of the classic mean-variance model,  \cite{Ortiz2023} employed covariance matrix shrinkage with ridge regression to stabilize portfolio weights and improve out-of-sample Sharpe ratios, thus making the model more practical. In the context of urban power decarbonization, \cite{Yang2024} integrated deep learning and green hydrogen storage into the Markowitz framework to enhance returns amid the challenges of renewable energy.\cite{Hui2024} proposed a Factor-Augmented Derating (FAD) technique for large-scale portfolios, narrowing the gap between expected and realized returns, with significant implications for stock market applications. \cite{Petukhina2024} further improved model stability by using projected gradient descent, which reduces sensitivity to input fluctuations and transaction costs, bolstering portfolio stability and out-of-sample performance.

Building on these core advancements, various innovative methods are being developed to refine stock selection, risk management, and overall portfolio performance. \cite{cheong2017using} introduced a two-stage approach combining k-means clustering for stock selection with genetic algorithms for optimizing portfolio weights. \cite{farid2022portfolio} employed self-organizing maps to cluster companies, followed by genetic algorithms to optimize selected stocks, creating a well-balanced portfolio. To address market instability, \cite{pahade2021multi} introduced a fuzzy logic model that considers variance and semi-variance, capturing market volatility for more robust optimization. Other structured search techniques have also emerged; for instance, \cite{akbay2020parallel} used neighborhood search algorithms with quadratic programming (QP) for optimal stock weights, while \cite{kizys2019biased} integrated assignment problem solutions with repetitive local search and QP for enhanced accuracy.

Sparse and short-term-focused models further contribute to diversified strategies. \cite{lai2018short} proposed a short-term, sparse model that simplifies portfolio composition by emphasizing the immediate profitability of selected stocks. Bayesian learning is another area of innovation, where \cite{de2019bayesian} used Bayesian algorithms with Gaussian assumptions to refine stock return estimates through prior and posterior distributions. For multi-period modeling,\cite{li2022multi} demonstrated that stepwise adjustments to portfolios over multiple periods could yield higher returns compared to single-period models. \cite{mercurio2020entropy} introduced an entropy minimization approach as an alternative to variance-based risk measures, creating a diversified portfolio structure that balances return and risk.

Advanced predictive and time-varying models have also gained traction. (\cite{katsikis2022time}) explored time-varying mean-variance datasets with time-variant quadratic programming (TVQP), while(\cite{kaczmarek2022building}) utilized machine learning to dynamically predict stock performance and optimize allocations. Meanwhile, (\cite{torrente2022rescaling}) applied game-theoretic adjustments to classical models, introducing competitive rebalancing for portfolio stability. (\cite{arcuri2023robust}) employed a Value at Risk (VaR) model that emphasizes strong conditional VaR to maintain stable, long-term performance.

Additional perspectives further enhance portfolio optimization, such as combining fundamental indexing (FI) with mean-variance optimization (\cite{pysarenko2019predictive}), and developing multivariate affine generalized hyperbolic (MAGH) frameworks (\cite{chaweewanchon2022markowitz}) for asset allocation. High-dimensional data in portfolio optimization is addressed by \cite{li2022spectrally}, who proposed a method for handling large, high-dimensional datasets effectively. \cite{katsikis2022time} introduced time-varying quadratic programming (TVQP) to manage datasets with changing means and variances over time.


This article introduce a novel concept called {\ttfamily"}well-behavedness of stocks{\ttfamily"} based on a retrospective idea.This approach predicts future stock performance by analyzing probability density functions of similar historical windows, thereby identifying stocks with consistent past behaviors likely to inform future trends. Ultimately, a suitable portfolio selection approach is combined with the main Markowitz approach. The algorithm's effectiveness is evaluated on the Tehran Stock Exchange and other stock exchanges worldwide. Section \ref{sec:Experiments} explains the proposed methods in detail in several subsections. In section \ref{sec:Discussion}, the authors introduce the dataset used, evaluation criteria, and the results of applying the proposed method and selecting effective portfolios. The final section concludes with a discussion and conclusion.

\section{Proposed Method}
\label{sec:Proposed}


The proposed method builds on the traditional Markowitz approach (\cite{markowitz1959portfolio}) and incorporates an adaptive framework for stock selection using elastic matching based on Kullback-Leibler (KL) divergence. The overall aim is to enhance portfolio optimization by identifying well-behaved stocks whose past behaviors can reliably predict future performance. The method is divided into several steps, as described below:


\subsection{Identifying Best Similar Window based on the function of similarity degree(FBS)}\label{sec:Find} 

Let the price of a stock over a time period $( T )$ be represented as a time series \( X_T = \{x_t\}_{t=1}^T \), where \( x_t \) is the price at time \( t \). The return signal \( r_t \) is derived from the price series, and it is used in further calculations for portfolio construction.

To predict future stock behavior, the proposed method searches for a similar window in the past. Assuming that the current window of interest starts at time \( t \) with a length of \( w \) days, the algorithm compares this window to past windows of the same stock. However, instead of searching the entire history, the search is limited to \( w'' \) days prior to \( t \).  the pseudo-code for the similarity search is provided in Algorithm \ref{alg:best-match}.

\begin{algorithm}
\caption{Finding the best matching window in stock history based on similarity degree}\label{alg:best-match}
\begin{algorithmic}
\State $ Best-similar-index \gets t$
\State $Best-divergence \gets 1000$
\For{$t^{\prime}=t+w:t+w^{\prime}$}
\State $divergence \gets D(r(t:t+w)\parallel (r:t^{\prime}:t^{\prime}+w))$
\If{$ divergence < Best-divergence$}
\State $Best-divergence \gets divergence$
\State $Best-similar-index \gets t^{\prime}$
\EndIf
\EndFor
\State \textbf{return} $Best-similar-index, Best-divergence$
\end{algorithmic}
\end{algorithm}
The divergence \( D \) between two windows is measured using kernel density estimates (KDE), as outlined below.

\subsection{Calculating the Degree of Divergence}
\label{sec:Function}

The divergence or dissimilarity between two windows is calculated using their probability density functions (PDFs). Instead of treating each window as a single unit, the method divides each \( w \)-day window into overlapping sub-windows of length \( w' \). 

For each sub-window pair \( i \), the divergence is calculated using a kernel density estimation (KDE), defined by :

\begin{equation} \label{kernel}
    \hat{f}_{t_i} = \frac{1}{n} \sum_{j=t+i}^{t+i+w'} K_h(x - r_j)
\end{equation}

Where \( K_h \) is the Gaussian kernel function. The total divergence is computed by averaging the divergences across all sub-windows by \ref{Div} :

\begin{equation} \label{Div}
    \text{div} = \sum_{i=0}^{w-w'}D(\hat{f}_{t_i} \parallel \hat{f}_{t'_i})
\end{equation}

This process is illustrated in Algorithm  \ref{alg:kde-divergence}, which provides the steps for calculating the dissimilarity between two windows.

\begin{algorithm}
\caption{Pseudocode for calculating the dissimilarity between two windows relative to each other}
\label{alg:kde-divergence}
\begin{algorithmic}
\State $div \gets 0$
\For{$i = 0$ \textbf{to} $w-w^{\prime}$}
\State $\text{Find}\; f1 = \text{KDE of}\; r(t+i:t+i+w^{\prime})$
\State $\text{Find}\; f2 = \text{KDE of}\; r(t^{\prime}+i:t^{\prime}+i+w^{\prime})$
\State $div \gets$ divergence of $f1$ and $f2$
\State $div \gets div + div$
\EndFor
\State \textbf{return} $div$
\end{algorithmic}
\end{algorithm}


where $\hat{f}_{t_i}$ and $\hat{f}_{t_j}$ are kernel density functions of two distributions  \ref{phat} and \ref{pprimehat} explain how these functions are calculated based on the Gaussian kernel:

\begin{equation} \label{phat}
    \hat{f}_{t_i}= \frac{1}{n(2\pi \sigma^2)^\frac{d}{2}}\sum_{i=1}^{n}exp(-\frac{\left\| x-x_{i} \right\|^2}{2\sigma ^2})
\end{equation}

\begin{equation} \label{pprimehat}
    \hat{f}_{t_j}= \frac{1}{n^{\prime}(2\pi {\sigma^{\prime}}^2)^\frac{d}{2}}\sum_{i^{\prime}=1}^{n^{\prime}}exp(-\frac{\left\| x-x^{\prime}_{i^{\prime}} \right\|^2}{2{\sigma^{\prime}}^2 })
\end{equation}

If the Jensen-Shannon divergence measure is used to calculate the distance between these two distributions, it should be calculated according to:
\begin{equation}\label{JS}
    D_{JS}(\hat{f}_{t_i}\parallel \hat{f}_{t_j})=\frac{1}{2}\left(D_{KL}\left[ \hat{f}_{t_i}\parallel\frac{\hat{f}_{t_i}+\hat{f}_{t_j}}{2} \right] + D_{KL}\left[ \hat{f}_{t_j}\parallel\frac{\hat{f}_{t_i}+\hat{f}_{t_j}}{2} \right]\right)
\end{equation}

KL stands for Kullback-Leibler divergence between two sub-windows \( t_i \) and \( t_j \) is defined as:

\begin{equation}\label{eq:kl_divergence_subwindows}
    D_{\text{KL}}(\hat{f}_{t_i} \parallel \hat{f}_{t_j}) = \sum_{x \in \mathcal{X}} \hat{f}_{t_i}(x) \log\left(\frac{\hat{f}_{t_i}(x)}{\hat{f}_{t_j}(x)}\right),
\end{equation}

where \( \mathcal{X} \) is the set of possible stock prices.


\subsection{Elastic Matching and Well-Behaved Stocks}
\label{sec:welbehaved}


The elastic matching condition evaluates the similarity between the recent behavior of a stock's price distribution and its past behavior by computing the mean Kullback-Leibler (KL) divergence. Let \( T_1 \) represent the most recent data window and \( T_i \) for \( i = 2, \dots, n \) represent past data windows. The mean KL divergence \( \overline{D_{KL}} \) between the current window \( T_1 \) and the past windows \( T_i \) is given by:

\begin{equation}\label{eq:D_KL}
    \overline{D_{KL}} = \frac{1}{n-1} \sum_{i=2}^{n} D_{KL}(F_{T_1} \parallel F_{T_i}).
\end{equation}

If the mean divergence \( \overline{D_{KL}} \) is smaller than a predefined threshold \( \epsilon \), i.e., \( \overline{D_{KL}} < \epsilon \), the stock's behavior over these sub-windows is considered elastic. This means there is a consistent pattern between the present and past distributions.

A stock is considered well-behaved if the elastic matching condition continues to hold over an extended period \( T + \Delta T \), where \( n' = \lfloor \Delta T / \tau \rfloor \) future sub-windows \( T_i \) for \( i = n+1, \dots, n+n' \) are examined. The stock's price distribution must satisfy:

\begin{equation}\label{eq:wel}
    \frac{1}{n'} \sum_{i=n+1}^{n+n'} D_{KL}(F_{T_1} \parallel F_{T_i}) \leq \epsilon.
\end{equation}

When this condition is met, the stock is interpreted as having predictable future behavior based on its past, indicating good performance over time.

\begin{algorithm}
\caption{The pseudocode to find the top 35 well-behaved stocks}
\label{alg:pseudocode}
\begin{algorithmic}
\State $\text{Indexs} \gets \{\}$
\For{$i=1$ \textbf{to} $N$}
\State $numberOfTrueCondition \gets 0$
\State $totalNumber \gets 0$
\For{$t=21$ \textbf{to} $\text{length}(r_i)$}
\State $\mu_{it} \gets \text{mean}(r_i(t-w:t))$
\State $\text{bestSimilarIndex}, \text{bestSimilarity} \gets \text{FBS}(r_i,t)$
\State $\mu_{it}' \gets \text{mean}(r_i(\text{bestSimilarIndex}-w:\text{bestSimilarIndex}))$
\If{$\text{bestSimilarity} < 1$}
\State $totalNumber \gets totalNumber + 1$
\If{$\mu_{it}/\mu_{it}' > 0$}
\State $numberOfTrueCondition \gets numberOfTrueCondition + 1$
\EndIf
\EndIf
\EndFor
\State $\text{Indexs(end+1)}.\text{index} \gets i$
\State $\text{Indexs(end)}.\text{percentage} \gets numberOfTrueCondition/totalNumber$
\EndFor
\State $\text{sortedIndex} \gets \text{sort}(\text{Indexs}, -\text{percentage})$
\State \textbf{return} $\text{sortedIndex}(1 : 35)$
\end{algorithmic}
\end{algorithm}

Algorithm\ref{alg:pseudocode} outlines the pseudocode used to identify well-behaved stocks. For each stock \( i \), we iterate over its price data \( r_i \), calculate the average return \( \mu_{it} \) over a window of length \( w \), and find the most similar past window using the FBS method. If the similarity score is below a threshold (typically \( <1 \)), and the ratio \( \mu_{it} / \mu'_{it} > 0 \) (where \( \mu'_{it} \) represents the average return of the similar window), the stock is classified as well-behaved for that window. The percentage of such well-behaved instances is computed over the entire time series, and the top 35 stocks with the highest percentage are selected for portfolio management.

\subsection{integrating Well-Behaved Stocks with the Markowitz's basic process}
\label{sec:combining}


Once the top well-behaved stocks are identified, the FBS method is integrated into the Markowitz optimization framework. In this approach, two data windows are considered: the current window and a future window identified as similar to the current one.As shown in Fig.\ref{fig:similarwindow}, the black window represents the current data window, while the green window represents the unknown future. The red window represents a past window that is similar to the current one, identified using the FBS algorithm. The purple window represents the future of this similar window. If the future of the similar window closely resembles the current window's future, this stock behavior is considered a good basis for portfolio decision-making.

\begin{figure}
\begin{center}
  \includegraphics[width=0.99\linewidth]{similarwindow.jpg}
  \caption{The meaning of current, future, similar, and future  of similar window in an example of implementing the FBS algorithm for a specific stock}
  \label{fig:similarwindow}  
\end{center}
\end{figure}

To quantify the similarity between the current and similar windows, we calculate the Jensen-Shannon (JS) divergence over 10 windows. If the total divergence is less than 1, the performance is typically better. A smaller divergence value indicates higher similarity, and if the divergence between the two windows is less than 0.5, more emphasis is placed on the future of the similar window for prediction.

The weight coefficient \( \gamma_i \) for the \( i \)-th stock in the portfolio is defined as follows:

\begin{equation}\label{gamma}
    \gamma_i = 1 - \min\left(1, \frac{0.5}{d}\right),
\end{equation}

where \( d \) is the divergence between the current and similar windows. This coefficient adjusts the importance placed on the future window versus the current window.

The Markowitz optimization model is then modified to account for both windows. The objective function becomes:


\begin{equation}\label{min}
\begin{split}
    min(& \lambda \left\lceil \sum_{i=1}^{N} \gamma_i \sum_{j=1}^{N}w_i w_j\sigma_{ij} \right\rceil
    -(1-\lambda) \left[ \sum_{i=1}^{N}\gamma_i w_i \mu_i \right]) \\ & + (\lambda^{\prime} \left\lceil \sum_{i=1}^{N} (1-\gamma_i) \sum_{j=1}^{N}w_i w_j \sigma^{\prime}_{ij} \right\rceil - (1-\lambda^{\prime}) \left[ \sum_{i=1}^{N}(1-\gamma_i) w_i \mu^{\prime}_i \right])
\end{split}
\end{equation}

where \( \lambda \) and \( \lambda' \) represent risk coefficients, \( \sigma_{ij} \) and \( \sigma'_{ij} \) are the covariances of the stock returns in the current and future windows, and \( \mu_i \) and \( \mu'_i \) are the mean returns for the current and future windows, respectively.

By solving this optimization problem using quadratic programming, the desired weight allocation for each stock in the portfolio is determined. The weight constraints ensure that each stock holds between 5\% and 40\% of the portfolio. This method ensures that the selected stocks not only show consistent behavior in the past but are also positioned to perform well in the future, improving the overall portfolio's stability and return potential.




\section{Experiments}\label{sec:Experiments}
In the proposed method, several key parameters are set to optimize performance. The window length \(w\) is set to 20 days for both the proposed method and the Markowitz process. The rolling window length \(w′\), used for comparing different windows, is set to 10 days, resulting in smaller 10-day windows within each 20-day window. Additionally, the parameter \(w′′\) determines the look-back period for identifying similar windows and is set to 500 days (roughly two years). To calculate the appropriate bandwidth for kernel density estimation (KDE), the improved Sheater and Jones method is applied to the daily return signals of each stock. This approach is consistently used for all stages involving the stock’s performance evaluation.

\subsection{Experimental data} \label{sec:Experimental data}
The experiments were conducted on daily stock price data from the Tehran Stock Exchange. Stocks that did not meet the behavior criteria, such as {\ttfamily"}Exchange-Traded Fund{\ttfamily"} and {\ttfamily"}Rights Offering Shares{\ttfamily"}, were excluded. Ultimately, 368 symbols were selected for testing the proposed algorithm. This selection ensures that only stocks exhibiting favorable characteristics for exploration are included in the analysis.


\subsection{Evaluation by Back Testing Method}
\label{sec:Eval}
To evaluate the performance of the proposed method, a back-testing system was implemented. Based on the system’s suggested coefficients for portfolio management at time \(t\), stocks were purchased. The algorithm measures the average daily return over an \(N\)-day window, using multiple-day windows for buying and selling decisions. At the end of each window, profits and losses are calculated, and the process is repeated over a total period of 1000 days. The evaluation is reported in ten 100-day intervals, with both short-term and overall average returns used to assess the algorithm’s effectiveness. A better performance over time indicates a more successful portfolio management strategy.

\subsection{FBS Method and Bandwidth Selection}
\label{sec:Eval FBS}

The FBS (Find Best Similarity) algorithm was tested to determine the optimal window size. Fig.\ref{fig:fbs} demonstrates the results for the “Foolad” stock symbol on day 537, where the red window represents the current time and the green window represents a similar past window. The future trends of both windows are closely aligned, showcasing the FBS algorithm's capability to identify key patterns. Additionally, bandwidth selection for KDE was tested using different methods. The improved Sheater and Jones method consistently outperformed other approaches. Fig.\ref{fig:comparing} illustrates how this method provided superior accuracy in predicting future trends, leading to its selection for the final approach.


\begin{figure}
\begin{center}
    \includegraphics{fbs.jpg}
    \caption{Executing the proposed FBS algorithm on the {\ttfamily"}Foolad{\ttfamily"} stock symbol, on the 537th day from the end}
    \label{fig:fbs}
\end{center}
\end{figure}




\begin{figure}
\begin{center}
  \includegraphics[width=0.9\linewidth]{comparing.jpg}
  \caption{Comparing different bandwidths for assessing the accuracy of the first condition using the best matching window when the similarity value exceeds a specified threshold}  
  \label{fig:comparing}
\end{center}
\end{figure}

\subsection{Identifying Well-Behaved Stocks}

The method defines well-behaved stocks based on ranking. Initially, stocks from the last 1000 days were analyzed, but to avoid bias from known back-testing results, this period was excluded from the final analysis. Instead, we focused on identifying well-behaved stocks from the preceding 2000 days, ensuring a forward-looking approach. Table \ref{table:finding35} presents the top 35 well-behaved symbols from the second thousand-day period, which were then used to continue the evaluation process for the first thousand days, maintaining the integrity of the test results.



\begin{longtable}{ p{.05\textwidth} p{.13\textwidth} p{.14\textwidth} p{.25\textwidth} p{.13\textwidth} p{.15\textwidth} } 
\caption{Lists the top 35 stocks identified using the Sheater and Jones bandwidth method, considering only the best similar windows that exceed a specified similarity threshold, while excluding the last thousand days of data.}
\label{table:finding35}
\\ \hline 
\textbf{Row} & \textbf{Symbol Stock Name} & \textbf{Accuracy Percentage}& \textbf{Industry} & \textbf{Total Number of Data Point} & \textbf{Number of instances where divergence exceeds the threshold} \\ \hline
\endhead
1 & Zegoldasht & 100 & Farm and related services & 125 & 1 \\ 
3 & Ghgolsta & 81.8182 & Food and Beverage Products & 125 & 9 \\ 
4 & Koravi & 81.8182 & Metal ore extraction & 444 & 9 \\ 
5 & Benirou & 76.4706 & Electrical machinery and equipment & 397 & 52 \\ 
6 & Setran & 74.1379 & Cement, lime, and gypsum & 454 & 43 \\ 
7 & Seydkoo & 72.0588 & Cement, lime, and gypsum & 166 & 49 \\ 
8 & VaAzar & 69.2308 & Construction, real estate, and properties & 363 & 54 \\ 
9 & Sheklar & 68.75 & Chemical products & 367 & 11 \\ 
10 & Kama & 68.2927 & Metal ore extraction & 376 & 28 \\ 
11 & Sekhrz & 68.1034 & Cement, lime, and gypsum & 379 & 79 \\ 
12 & FarAvar & 67.8571 & Basic metals & 383 & 19 \\ 
13 & Sepaha & 67.3469 & Cement, lime, and gypsum & 436 & 33 \\ 
14 & Chekaveh & 66.6667 & Paper products & 400 & 32 \\ 
15 & Shenfat & 65.6051 & Petroleum, coke, and nuclear fuel products & 451 & 103 \\ 
16 & Ghadasht & 65.285 & Food and Beverage Products excluding sugar and sweets & 376 & 126 \\ 
17 & Nbours & 64.7059 & Auxiliary activities to financial intermediaries & 117 & 11 \\ 
18 & Kenour & 63.2911 & Metal ore extraction & 177 & 50 \\ 
19 & Barkat & 62.5 & Pharmaceutical materials and products & 114 & 10 \\ 
20 & Kesaveh & 62.2642 & Tiles and ceramics & 372 & 33 \\ 
21 & Kesaedi & 61.8421 & Tiles and ceramics & 360 & 47 \\ 
22 & Deshimi & 61.7647 & Pharmaceutical materials and products & 377 & 84 \\ 
23 & Faspa & 61.5 & Basic metals & 362 & 123 \\ 
24 & Fasmin & 61.2903 & Basic metals & 410 & 19 \\ 
25 & Kebourse & 61.1111 & Auxiliary activities to financial intermediaries & 117 & 11 \\ 
26 & Bourse & 61.1111 & Auxiliary activities to financial intermediaries & 117 & 11 \\ 
27 & Ghapira & 61.039 & Sugar and sweets & 375 & 47 \\ 
28 & Maroun2 & 60.5839 & Chemical products & 238 & 83 \\ 
29 & Sarbil & 60.2804 & Cement, lime, and gypsum & 325 & 129 \\ 
30 & Maroun & 59.4203 & Chemical products & 238 & 82 \\ 
31 & Fakhas & 59.0517 & Basic metals & 331 & 137 \\ 
32 & Sashahed & 58.9744 & Construction,real estate, and properties & 365 & 23 \\ 
33 & Depars & 58.9286 & Pharmaceutical materials and products & 458 & 99 \\ 
34 & Safha & 58.7065 & Cement, lime, and gypsum & 346 & 118 \\ 
35 & Fama & 58.6319 & Metal products manufacturing & 405 & 180 \\ \hline
% needs to go inside longtable environment
\label{tab:myfirstlongtable}
\end{longtable}


\subsection{Comparison of the Proposed Method and the Markowitz Baseline Method}
In this section, we utilize the Markowitz baseline method to calculate the mean return rate for the 50 stock portfolios obtained in the previous section, employing a back-testing approach. We repeat this process for the proposed method and then compare the results from both methods. By subtracting the daily average return of the Markowitz method from that of the proposed method, we can draw conclusions regarding the superiority of one approach over the other and assess the degree of improvement achieved as shown in Table \ref{table:improvement}.

\begin{longtable}{ p{.15\textwidth}  p{.27\textwidth} p{.25\textwidth} p{.25\textwidth} }
\caption{The percentage of improvement in the proposed method compared to the Markowitz baseline method for 50 different portfolios.}
\label{table:improvement}
% \begin{tabularx}{\textwidth}{|c|X|X|X|}
\hline 
\textbf{Portfolio} & \textbf{The result of subtracting B from A, divided by the absolute value of B, multiplied by 100} & \textbf{Average return using the Markowitz baseline method (B)} & \textbf{Average return using the proposed method (A)} \\ 
\hline
\endhead
1& 20.7415 & 0.2573 & 0.3106 \\
2 & 56.2041 & 0.1801 & 0.2814 \\
3 & 45.4560 & 0.2566& 0.3732 \\
4 & 152.7466& 0.1062& 0.2683 \\
5& 15.6404& 0.1904& 0.2202 \\
6& 68.2275& 0.1604& 0.2698 \\
7 & 57.6507& 0.1643& 0.2591 \\
8 & 63.2526& 0.2035& 0.3322                    \\
9& 1.4891& 0.2911& 0.2955                    \\
10& 24.5071& 0.2753& 0.3428                    \\
11 & 12.344 & 0.2757 & 0.2417                    \\
12 & 104.9510 & 0.1340& 0.2747                    \\
13 & 57.7825& 0.1292& 0.2039                    \\
14 & 8.2209& 0.3380& 0.3103                    \\
15 & 2.3804 & 0.2232& 0.2285                    \\
16 & 9.7994 & 0.2241& 0.2460                    \\
17 & 38.6751 & 0.1954& 0.2709                    \\
18 & 7.1199 & 0.2488 & 0.2665                    \\
19 & 45.8413& 0.1851 & 0.2700                    \\
20 & 1.9877& 0.2048 & 0.2088                    \\
21 & 40.2888 & 0.2051 & 0.2878                    \\
22 & 17.1097  & 0.1616 & 0.1893                    \\
23 & 50.8998 & 0.2845  & 0.4293                    \\
24 & 16.3096  & 0.2214 & 0.2575                    \\
25 & 42.4601  & 0.2629 & 0.1513                    \\
26 & 26.5798  & 0.2824  & 0.3575                    \\
27 & 15.8851 & 0.2045 & 0.2369                    \\
28 & 46.1852 & 0.1982 & 0.2897  \\
29 & 49.9436 & 0.2034 & 0.3050                    \\
30 & 3.7192 & 0.2909 & 0.2801                    \\
31 & 4.0503 & 0.2324  & 0.2230                    \\
32 & 2.3509 & 0.2526 & 0.2585                    \\
33 & 34.7096  & 0.1422 & 0.1915                    \\
34 & 9.3057  & 0.2628  & 0.2873                    \\
35 & 8.2149  & 0.2700 & 0.2922                    \\
36 & 44.9178 & 0.2109 & 0.3057                    \\
37 & 54.6303 & 0.2706 & 0.4184                    \\
38 & 17.3260 & 0.2983 & 0.3500                    \\
39 & 19.3977 & 0.2257 & 0.2695                    \\
40 & 9.2423 & 0.1750 & 0.1912                    \\
41 & 34.1795 & 0.2613 & 0.3505                    \\
42 & 39.0505 & 0.2474 & 0.3440                    \\
43 & 26.8253 & 0.2657 & 0.1944                    \\
44 & 15.5447 & 0.1981 & 0.1673                    \\
45 & 127.3386  & 0.1286 & 0.2923                    \\
46 & 45.3645 & 0.1488 & 0.2163                    \\
47 & 12.2787 & 0.2471 & 0.2774                    \\
48  & 22.2219 & 0.2218 & 0.2711                    \\
49 & 47.8972 & 0.2152 & 0.3183                    \\
50 & 32.5290 & 0.2455 & 0.3254                    \\ \hline
% \end{tabularx}
\noalign{\smallskip}
\end{longtable}


\section{Discussion and Conclusion}\label{sec:Discussion}
Table \ref{table:comparing}  the performance improvements of the proposed method compared to the baseline Markowitz approach across fifty portfolios. The behavior of the portfolios was evaluated during the second thousand days, while the back-testing was conducted over the most recent thousand days. The results clearly highlight the benefits of the proposed method in identifying well-behaved stocks and leveraging future windows to enhance performance.

\begin{table}
\caption{Comparison of the proposed method with the Markowitz baseline for 50 portfolios of well-behaved stocks over the last 1000 days.}
\label{table:comparing}
\begin{tabular}{c  p{0.7\linewidth}}
  \hline
  \textbf{Proposed method} & \multicolumn{1}{c}{\textbf{Results}} \\ 
  \hline
  7\% & Percentage of portfolios with negative improvement, indicating worse performance compared to the Markowitz baseline.  \\ 
  
  86\% & Success rate of the proposed method, calculated as the number of portfolios with improved performance divided by the total number of portfolios. $(\frac{43}{50}*100)$ \\ 
  
  29.7\% & Average percentage improvement across all cases. \\
  
  37.2\% & Average percentage improvement in cases where the method led to better performance. \\
  \hline
\end{tabular}
\label{tab:compared}
\end{table}

The result shows that, when well-behaved stocks are accurately identified and the proposed method is applied, significant improvements in portfolio performance can be achieved. This reinforces the value of integrating the current and future windows in the optimization process for more accurate forecasting and enhanced portfolio returns.


Future work can improve upon the current methodology by introducing intelligent stock selection criteria, such as using fundamental stock indicators or clustering based on recent performance, rather than random selection. Additionally, the Markowitz optimization process could be enhanced by incorporating variable or multi-length windows and adjusting the risk-taking coefficient \( \lambda \) over time to better capture market dynamics.

Other potential enhancements include refining the back-testing process by accounting for trading factors like buy and sell queues, which would more accurately reflect real-world conditions. Improvements to the FBS algorithm, such as better distance learning between similar future windows and their actual outcomes, along with the use of variable bandwidth over time, could also lead to more precise and adaptive forecasting models.

%% If you have bibdatabase file and want bibtex to generate the
%% bibitems, please use
%%
 %\bibliographystyle{plainnat} 
 %\bibliographystyle{agsm} % Name-Date (Harvard) style
 %\bibliography{cas-refs}



\bibliographystyle{elsarticle-harv} 
\bibliography{cas-refs}

\end{document}
\endinput
